\section{Motivation and Questions}

\begin{frame}{Motivation}
    \begin{itemize}
        \item Attention-based models dominate language, vision, and multimodal benchmarks, but we still lack reliable local explanations for specific outputs.
        \item This proposal focuses on three connected fronts:
        \begin{itemize}
            \item emergence during training (grokking-like transitions),
            \item controllable neuro-symbolic steering,
            \item shift/attack diagnostics from internal features.
        \end{itemize}
        \item Core question: can we connect internal mechanisms to measurable, controllable behavior with explicit trade-offs?
    \end{itemize}
    \vspace{0.3em}
    \footnotesize References: \cite{vaswani2017attention,brown2020gpt3,openai2023gpt4,kautz2022thirdaisummer}
\end{frame}

\begin{frame}{Motivation: Why Local Explanations Matter}
    \small
    \begin{itemize}
        \item Current models perform well, but output quality alone does not explain \textit{why} a specific decision was produced.
        \item Without local explanations, we cannot reliably separate:
        \begin{itemize}
            \item real reasoning from spurious shortcuts,
            \item robust behavior from fragile behavior under shift,
            \item targeted control from collateral regressions.
        \end{itemize}
        \item The proposal therefore prioritizes measurable internal evidence, not only end-task scores.
    \end{itemize}
\end{frame}

\begin{frame}{Motivation: Emergence (G1)}
    \small
    \begin{block}{Question}
        How do useful internal structures emerge during training, especially around delayed generalization (grokking-like transitions)?
    \end{block}
    
    \begin{itemize}
        \item Analyze phase transitions in train/test behavior.
        \item Measure feature stability and cross-layer consistency around transition windows.
        \item Use controlled distribution shifts to stress-test emergence conditions.
    \end{itemize}
\end{frame}

\begin{frame}{Motivation: Neuro-symbolic Steering (G2)}
    \small
    \begin{block}{Question}
        Can we steer targeted behavior with logic constraints while keeping quality loss and collateral effects bounded?
    \end{block}
    \begin{itemize}
        \item Use sparse features and LTN validators as control signals.
        \item Compare constrained decoding against unconstrained and lightweight baselines.
        \item Report explicit quality-versus-satisfaction trade-offs.
    \end{itemize}
\end{frame}

\begin{frame}{Motivation: Diagnostics (G3)}
    \small
    \begin{block}{Question}
        Are internal sparse signatures better early-warning signals for shift/attacks than output-only heuristics?
    \end{block}
    \begin{itemize}
        \item Evaluate clean vs perturbed separation in latent space.
        \item Measure ROC-AUC, PR-AUC, calibration, and lead-time.
        \item Test transfer to settings where output confidence can be misleading.
    \end{itemize}
\end{frame}

\begin{frame}{Motivation Detail: Core Question in Operational Form}
    \small
    \begin{block}{Core question}
        Can we connect internal mechanisms to controllable behavior with explicit and reproducible trade-offs?
    \end{block}
    \vspace{0.4em}
    \begin{table}
        \centering
        \begin{tabular}{p{0.28\linewidth}p{0.62\linewidth}}
            \toprule
            \textbf{If yes, we should observe} & \textbf{How we verify} \\
            \midrule
            Stable local features & Cross-seed/epoch stability and reconstruction fidelity \\
            Targeted control & Constraint gain with bounded quality/regression cost \\
            Reliable diagnostics & Internal-signal performance at least matching output-only baselines \\
            \bottomrule
        \end{tabular}
    \end{table}
\end{frame}

\begin{frame}{Core Claims and Research Questions}
    \small
    \begin{table}
        \centering
        \begin{tabular}{p{0.18\linewidth}p{0.35\linewidth}p{0.38\linewidth}}
            \toprule
            \textbf{Claim} & \textbf{Research question} & \textbf{Hypothesis summary} \\
            \midrule
            C1: Emergence &
            Can controlled shifts induce delayed generalization with stable internal changes? &
            Late test gains align with increased cross-seed/cross-layer feature consistency. \\
            \addlinespace[0.3em]
            C2: Steering locality &
            Can logic constraints change targeted behavior with bounded regressions? &
            LTN-guided constraints improve rule satisfaction with acceptable quality loss. \\
            \addlinespace[0.3em]
            C3: Diagnostics &
            Are sparse internal signatures stronger than output-only heuristics for perturbation detection? &
            Sparse latent signatures improve or match detection/calibration stability. \\
            \bottomrule
        \end{tabular}
    \end{table}
\end{frame}

\begin{frame}{Proposal Positioning}
    \begin{columns}[T,onlytextwidth]
        \column{0.33\textwidth}
        \begin{block}{Completed}
            \footnotesize
            \begin{itemize}
                \item Controlled grokking evidence under distribution shift.
                \item Sparse SAE diagnostics for clean-vs-perturbed separation.
            \end{itemize}
        \end{block}
        \column{0.33\textwidth}
        \begin{block}{In Progress}
            \footnotesize
            \begin{itemize}
                \item Frozen steering protocol calibration.
                \item Final baseline and ablation matrix definition.
            \end{itemize}
        \end{block}
        \column{0.33\textwidth}
        \begin{block}{Pending}
            \footnotesize
            \begin{itemize}
                \item Final multi-seed runs and transfer/failure analysis.
                \item Thesis integration and defense-ready narrative.
            \end{itemize}
        \end{block}
    \end{columns}
\end{frame}
